\documentclass[10pt]{article}
\usepackage[T1]{fontenc}
\usepackage[utf8]{inputenc}
\usepackage[french]{babel}
\usepackage{amsmath,amssymb,mathrsfs,mathpazo}
\usepackage{geometry}
\geometry{a4paper,margin=2cm}
\usepackage{graphicx,array,tabularx,tabularray,colortbl,multirow,multicol,booktabs,enumitem,pifont,ulem,fancybox,diagbox,cancel}
\usepackage{pstricks,pstricks-add,pst-plot,pst-tree,pst-node,pst-text,pst-func,pst-math,pst-bspline,pst-3dplot,pst-eucl,pst-solides3d}
\let\pscalculate\undefined
\usepackage{pst-bezier}
\usepackage{tikz,pgfplots,tkz-tab}
\usetikzlibrary{arrows,arrows.meta,patterns,calc,positioning,trees,decorations.markings,shapes.geometric,shapes.misc,shapes.arrows,intersections,decorations.pathmorphing,decorations.pathreplacing,backgrounds,shadows,babel}
\pgfplotsset{compat=1.18}
\usepgfplotslibrary{fillbetween}
\usepackage[np]{numprint}
\usepackage{listings,makecell,pas-tableur,twoopt}
\let\barre\undefined
\usepackage[pstricks]{tablvar}
\definecolor{darkgreen}{RGB}{0,100,0}
\providecommand{\red}{\color{red}}
\providecommand{\blue}{\color{blue}}
\providecommand{\green}{\color{green}}
\providecommand{\black}{\color{black}}
\providecommand{\white}{\color{white}}
\providecommand{\gray}{\color{gray}}
\providecommand{\cyan}{\color{cyan}}
\providecommand{\magenta}{\color{magenta}}
\providecommand{\yellow}{\color{yellow}}
\providecommand{\violet}{\color{violet}}
\providecommand{\orange}{\color{orange}}
\providecommand{\euro}{\texteuro}
\providecommand{\EUR}{\texteuro}
\providecommand{\ieme}{\textsuperscript{e}}
\providecommand{\degres}{\ensuremath{^\circ}}
\providecommand{\degree}{\ensuremath{^\circ}}
\providecommand{\vv}[1]{\overrightarrow{#1}}
\providecommand{\Coord}[1]{#1}
\providecommand{\decosix}{}
\providecommand{\starredbullet}{\ensuremath{\bullet}}
\pagestyle{empty}
\setlength{\parindent}{0pt}
\usepackage{adjustbox}\DeclareUnicodeCharacter{2236}{:}\DeclareUnicodeCharacter{2190}{\ensuremath{\leftarrow}}
\providecommand{\R}{\mathbb{R}}
\providecommand{\N}{\mathbb{N}}
\providecommand{\Z}{\mathbb{Z}}
\providecommand{\Q}{\mathbb{Q}}
\providecommand{\D}{\mathbb{D}}
\providecommand{\C}{\mathbb{C}}
\providecommand{\e}{{\mathrm{e}}}
\providecommand{\vect}[1]{\overrightarrow{#1}}
\providecommand{\vectt}[1]{\overrightarrow{\mathrm{#1}}}
\providecommand{\vvt}[1]{\overrightarrow{\mathrm{#1}}}
\providecommand{\np}[2][]{{#2~#1}}
\providecommand{\vv}[1]{{\overrightarrow{#1}}}
\renewcommand{\d}{{\mathrm{d}}}
\providecommand{\degres}{{^\circ}}
\providecommand{\degre}{{^\circ}}
\providecommand{\blue}{\color{blue}}
\providecommand{\red}{\color{red}}
\providecommand{\enskip}{\quad}
\providecommand{\llbracket}{\lbrack\!\lbrack}
\providecommand{\rrbracket}{\rbrack\!\rbrack}
\providecommand{\euro}{{\text{€}}}
\providecommand{\no}{{\text{n°}}}
\providecommand{\interval}[2]{{[#1;#2]}}
\providecommand{\intervFO}[2]{{[#1;#2[}}
\providecommand{\dotfill}{\ldots}

\providecommand{\ds}{\displaystyle}
\providecommand{\textding}[1]{#1}
\providecommand{\Oij}{\(\mathrm{O};\vec{\imath},\vec{\jmath}\)}
\providecommand{\Oijk}{\(\mathrm{O};\vec{\imath},\vec{\jmath},\vec{k}\)}
\providecommand{\Ouv}{\(\mathrm{O};\vec{u},\vec{v}\)}
\providecommand{\intervFF}[2]{\ensuremath{[#1;#2]}}
\renewcommand{\euro}{\ifmmode\text{\texteuro}\else\texteuro\fi}
\DeclareUnicodeCharacter{00B0}{\ifmmode^{\circ}\else\textdegree\fi}
\DeclareUnicodeCharacter{20AC}{\ifmmode\text{\texteuro}\else\texteuro\fi}
% Énoncés : Éducation nationale. Transcription et corrigés : APMEP, Denis Vergès et contributeurs. Découpage et classement : CoopMaths, Sébastien Lozano, Éric Elter et contributeurs. Adaptation de la présentation : Maths Explorer.
% CC BY-SA : https://coopmaths.fr/www/about/
% https://coopmaths.fr/alea/static/bac/2023/tex/bac_2023_09_sujet1_polynesie_4_cor.tex
\begin{document}
\noindent\textit{APMEP / CoopMaths. CC BY-SA. Présentation adaptée par Maths Explorer.}\par\medskip
\medskip





L'espace est muni d'un repère orthonormé \Oijk{} dans lequel on considère :

\begin{itemize}
\item[$\bullet$~~] les points A$(6~;~-6~;~6)$, B$(-6~;~0~;~6)$ et C$(-2~;~-2~;~11)$.
\item[$\bullet$~~] la droite $(d)$ orthogonale aux deux droites sécantes (AB) et (BC) et passant par le point A ;
\item[$\bullet$~~] la droite $\left(d'\right)$ de représentation paramétrique :

\[\left\{\begin{array}{l !{=} r l}
x &  -6 - 8t\\
y & 4t & \text{avec }\:t \in \R.\\
z &  6 + 5t
\end{array}\right.\]
\end{itemize}

\medskip

\textbf{Question 1}

Parmi les vecteurs suivants, lequel est un vecteur directeur de la droite $(d)$ ?

\begin{center}
\begin{tabularx}{\linewidth}{*{4}{X}}
\textbf{a.~} $\vect{u_1}\begin{pmatrix}-6\\3\\0\end{pmatrix}$&
\textbf{b.~} $\vect{u_2}\begin{pmatrix}1\\2\\6\end{pmatrix}$&
\textbf{c.~} $\vect{u_3}\begin{pmatrix}1\\2\\0,2\end{pmatrix}$&
\textbf{d.~} $\vect{u_4}\begin{pmatrix}1\\2\\0\end{pmatrix}$
\end{tabularx}
\end{center}

\medskip

\begin{tabular}{@{\hspace{0.05\linewidth}} | p{0.92\linewidth}}
Il faut déterminer quel vecteur est orthogonal à la fois à $\vectt{AB}$ et à $\vectt{BC}$.

$\vectt{AB} \begin{pmatrix} -12\\6\\0 \end{pmatrix}$ et
$\vectt{BC} \begin{pmatrix} 4\\-2\\5 \end{pmatrix}$

$\vect{u_2}\cdot\vectt{AB}=1\times (-12)+2\times 6 + 0\times 0=0$ donc $\vect{u_2}\perp \vectt{AB}$

$\vect{u_2}\cdot\vectt{BC}=1\times 4+2\times (-2) + 0\times 5=0$ donc $\vect{u_2}\perp \vectt{BC}$

\hfill\textbf{Réponse b.}
\end{tabular}

\medskip

\textbf{Question 2}

Parmi les équations suivantes, laquelle est une représentation paramétrique de la droite (AB) ?

\begin{center}
\begin{tabularx}{\linewidth}{*{2}{X}}
\textbf{a.~}$\left\{\begin{array}{l !{=} r}
x & 2t + 6\\
y &-6\\
z& t+6
\end{array}\right. \:\text{avec }\: t \in \R$&
\textbf{b.~}$\left\{\begin{array}{l !{=} r}
x & 2t - 6\\
y &- 6\\
z&- t - 6
\end{array}\right. \:\text{avec }\: t \in \R$\\
\textbf{c.~}$\left\{\begin{array}{l !{=} r}
x & 2t + 6\\
y &-t -6\\
z& 6
\end{array}\right. \:\text{avec }\: t \in \R$&
\textbf{d.~}$\left\{\begin{array}{l !{=} r}
x & 2t + 6\\
y & t - 6\\
z& 6
\end{array}\right. \:\text{avec }\: t \in \R$
\end{tabularx}
\end{center}

\medskip

\begin{tabular}{@{\hspace{0.05\linewidth}} | p{0.92\linewidth}}
La droite (AB) a pour vecteur directeur $\vectt{AB}$ et parmi les quatre propositions, le seul vecteur colinéaire à $\vectt{AB}$ est celui de coordonnées $(2\;;\;-1\;;\;0)$ correspondant à la réponse c.

\hfill\textbf{Réponse c.}
\end{tabular}


\textbf{Question 3}

Un vecteur directeur de la droite $\left(d'\right)$ est :

\begin{center}
\begin{tabularx}{\linewidth}{*{4}{X}}
\textbf{a.~} $\vect{v_1}\begin{pmatrix}-6\\0\\6\end{pmatrix}$&
\textbf{b.~} $\vect{v_2}\begin{pmatrix}-14\\4\\11\end{pmatrix}$&
\textbf{c.~} $\vect{v_3}\begin{pmatrix}8\\-4\\-5\end{pmatrix}$&
\textbf{d.~} $\vect{v_4}\begin{pmatrix}8\\4\\5\end{pmatrix}$
\end{tabularx}
\end{center}

\medskip

\begin{tabular}{@{\hspace{0.05\linewidth}} | p{0.92\linewidth}}
Un vecteur directeur de la droite $\left(d'\right)$ est 
$\vect{v} \begin{pmatrix} -8\\4\\5 \end{pmatrix}$ qui est l'opposé du vecteur $\vect{v_3}$.

\hfill\textbf{Réponse c.}
\end{tabular}

\medskip

\textbf{Question 4}

Lequel des quatre points suivants appartient à la droite $\left(d'\right)$ ?

\begin{center}
\begin{tabularx}{\linewidth}{*{2}{X}}
\textbf{a.~}M$_1(50~;~-28~;~-29)$&\textbf{b.~}M$_2(-14~;~-4~;~1)$\\
\textbf{c.~}M$_3(2~;~-4~;~-1)$& \textbf{d.~}M$_4 (-3~;~0~;~3)$
\end{tabularx}
\end{center}

\medskip

\begin{tabular}{@{\hspace{0.05\linewidth}} | p{0.92\linewidth}}
Il faut déterminer lequel des points M a des coordonnées qui vérifient le système
$\left \{
\begin{array}{l !{=} r}
x_{\text{M}} & -6-8t\\
y_{\text{M}} & 4t\\
z_{\text{M}} & 6+5t
\end{array}
\right .$

Pour M$_1$ on a: $\left \{
\begin{array}{r !{=} r}
50 & -6-8t\\
-28 & 4t\\
-29 & 6+5t
\end{array}
\right .$
qui est vérifié pour $t=-7$.

\hfill\textbf{Réponse a.}
\end{tabular}

\medskip

\textbf{Question 5}

Le plan d'équation $x = 1$ a pour vecteur normal:

\begin{center}
\begin{tabularx}{\linewidth}{*{4}{X}}
\textbf{a.~} $\vect{n_1}\begin{pmatrix}1\\0\\0\end{pmatrix}$&
\textbf{b.~} $\vect{n_2}\begin{pmatrix}0\\1\\1\end{pmatrix}$&
\textbf{c.~} $\vect{n_3}\begin{pmatrix}0\\1\\0\end{pmatrix}$&
\textbf{d.~} $\vect{n_4}\begin{pmatrix}1\\0\\1\end{pmatrix}$
\end{tabularx}
\end{center}

\medskip

\begin{tabular}{@{\hspace{0.05\linewidth}} | p{0.92\linewidth}}
Le plan d'équation $ax+by+cz+d=0$ a pour vecteur normal 
$\vect{n} \begin{pmatrix}a\\b\\c\end{pmatrix}$ donc le plan d'équation $x=1$ a pour vecteur normal le vecteur de coordonnées $\begin{pmatrix}1\\0\\0\end{pmatrix}$.

\hfill\textbf{Réponse a.}
\end{tabular}

\end{document}
